\section{Conclusão}

Nesse relatório, foi feita uma extensa revisão bibliográfica do estado da arte, cobrindo 3 tópicos: estimação do campo de velocidade do som (SSF) utilizando \textsl{Travel-Time Tomography} (TTT); inversão de SSF utilizando informação de tempo de trânsito; e inversão do SSF na coluna d'água. Também foi coberto, de maneira introdutória, alguns princípios ligados aos efeitos de parâmetros da água (como temperatura, pressão, e salinidade) no SSF, a relação do SSF com o comportamento de ondas sonoras na água, e algumas limitações relacionadas ao projeto.

A partir dessas revisões bibliográficas, foi possível estabelecer os fundamentos da TTT, e de como ela funciona, além de destacar diferenças entre suas implementações na literatura. Também foram selecionados artigos de relevância nas revisões de inversão usando tempo de trânsito, e inversão na coluna d'água, que foram posteriormente categorizados seguindo diversas categorias, buscando elencar os artigos que mais se encaixam no escopo e nas limitações do projeto.

Com isso, foram selecionadas 4 abordagens viáveis, (TTT, funções empíricas, parametrização, e redes neurais) a serem exploradas e adaptadas matematicamente ao problema sendo tratado no projeto, implementadas virtualmente e, se obtiverem resultados interessantes, implementadas com dados práticos.